\begin{frame}{Proof by Contradiction}
\begin{center}
\vspace{1.5cm}
{\Huge\textbf{Proof by Contradiction}}
\end{center}
\end{frame}

% --------------------------------------------------
% Question 1
% --------------------------------------------------

\begin{frame}{Question 1}
A scheduling algorithm claims that there is no smallest positive integer.

Determine whether this claim is correct. Use proof by contradiction to
justify your conclusion.
\end{frame}

\begin{frame}{Question 1 -- Answer}
The claim is false because $1$ is the smallest positive integer.

Assume, for contradiction, that there is no smallest positive integer.

Then there must be some positive integer $n$ that is smaller than every
positive integer.
\end{frame}

\begin{frame}{Question 1 -- Answer}
Since $n$ is positive:

$$
n\geq1
$$

But $1$ is also a positive integer.

If $n>1$, then:

$$
1<n
$$

which contradicts the assumption that $n$ is smaller than every positive
integer.
\end{frame}

\begin{frame}{Question 1 -- Answer}
Therefore, the assumption leads to a contradiction.

Hence:

$$
\boxed{1\text{ is the smallest positive integer}}
$$

and the original claim is false.
\end{frame}

% --------------------------------------------------
% Question 2
% --------------------------------------------------

\begin{frame}{Question 2}
A database system uses integer identifiers. A developer claims that there
can be two consecutive integers that are both even.

Prove that this is impossible using contradiction.
\end{frame}

\begin{frame}{Question 2 -- Answer}
Assume, for contradiction, that two consecutive integers $n$ and $n+1$
are both even.

Then there exist integers $a$ and $b$ such that:

$$
n=2a
$$

and:

$$
n+1=2b
$$

\end{frame}

\begin{frame}{Question 2 -- Answer}
Subtract the two equations:

$$
(n+1)-n=2b-2a
$$

Therefore:

$$
1=2(b-a)
$$

The right-hand side is even, while the left-hand side is odd.

This is a contradiction.
\end{frame}

\begin{frame}{Question 2 -- Answer}
Therefore, two consecutive integers cannot both be even.

Hence:

$$
\boxed{\text{No two consecutive integers are both even.}}
$$

\end{frame}

% --------------------------------------------------
% Question 3
% --------------------------------------------------

\begin{frame}{Question 3}
A software system requires every integer value in a particular set to be
both even and odd.

Show that such an integer cannot exist.
\end{frame}

\begin{frame}{Question 3 -- Answer}
Assume, for contradiction, that an integer $n$ is both even and odd.

Since $n$ is even:

$$
n=2a
$$

for some integer $a$.
\end{frame}

\begin{frame}{Question 3 -- Answer}
Since $n$ is also odd:

$$
n=2b+1
$$

for some integer $b$.

Therefore:

$$
2a=2b+1
$$

which gives:

$$
2(a-b)=1
$$

\end{frame}

\begin{frame}{Question 3 -- Answer}
The left-hand side is even, but the right-hand side is odd.

This is impossible.

Therefore:

$$
\boxed{\text{No integer is both even and odd.}}
$$

\end{frame}

% --------------------------------------------------
% Question 4
% --------------------------------------------------

\begin{frame}{Question 4}
A network engineer claims that the sum of two odd integers can be odd.

Determine whether the claim is possible. Prove your conclusion by
contradiction.
\end{frame}

\begin{frame}{Question 4 -- Answer}
Assume, for contradiction, that two odd integers $a$ and $b$ have an odd
sum.

Since both are odd:

$$
a=2m+1
$$

and:

$$
b=2n+1
$$

\end{frame}

\begin{frame}{Question 4 -- Answer}
Their sum is:

$$
a+b=2m+1+2n+1
$$

Therefore:

$$
a+b=2(m+n+1)
$$

Thus, $a+b$ is even.

This contradicts the assumption that the sum is odd.
\end{frame}

\begin{frame}{Question 4 -- Answer}
Hence:

$$
\boxed{\text{The sum of two odd integers is always even.}}
$$

\end{frame}

% --------------------------------------------------
% Question 5
% --------------------------------------------------

\begin{frame}{Question 5}
A computing algorithm claims that the square root of $2$ can be represented
exactly as a rational number.

Prove that this is impossible using contradiction.
\end{frame}

\begin{frame}{Question 5 -- Answer}
Assume, for contradiction, that:

$$
\sqrt{2}=\frac{p}{q}
$$

where $p$ and $q$ are integers with no common factor.

Squaring both sides:

$$
2=\frac{p^2}{q^2}
$$

\end{frame}

\begin{frame}{Question 5 -- Answer}
Therefore:

$$
p^2=2q^2
$$

Hence $p^2$ is even.

Therefore $p$ must be even.

Let:

$$
p=2k
$$

\end{frame}

\begin{frame}{Question 5 -- Answer}
Substitute into:

$$
p^2=2q^2
$$

giving:

$$
(2k)^2=2q^2
$$

$$
4k^2=2q^2
$$

Therefore:

$$
q^2=2k^2
$$

So $q$ is also even.
\end{frame}

\begin{frame}{Question 5 -- Answer}
Thus both $p$ and $q$ are even.

This contradicts the assumption that $p$ and $q$ have no common factor.

Therefore:

$$
\boxed{\sqrt{2}\text{ is irrational}}
$$

\end{frame}
